\documentclass[11pt,a4paper]{article}

\usepackage[polish]{babel}
\usepackage[T1]{fontenc}
\usepackage[utf8]{inputenc}
\usepackage{lmodern}
\usepackage[margin=2.5cm]{geometry}
\usepackage{tcolorbox}
\usepackage{enumitem}
\usepackage[hidelinks]{hyperref}
\usepackage{titlesec}
\usepackage{xcolor}
\usepackage{parskip}
\usepackage{amsmath}

\definecolor{granatowy}{HTML}{1a1a2e}
\definecolor{ciemnyniebieski}{HTML}{16213e}
\definecolor{jasnyniebieski}{HTML}{e7f3ff}

\titleformat{\section}{\Large\bfseries\color{granatowy}}{\thesection.}{0.5em}{}
\titleformat{\subsection}{\large\bfseries\color{ciemnyniebieski}}{\thesubsection}{0.5em}{}

\tcbuselibrary{skins,breakable}
\newtcolorbox{infobox}{
	colback=jasnyniebieski, colframe=jasnyniebieski, sharp corners, boxrule=0pt,
	left=10pt, right=10pt, top=8pt, bottom=8pt, breakable
}

\title{\textbf{Skąd biorą się liczby na mapie?}\\
	\large Przystępne wyjaśnienie mapy potencjału lokalizacji\\
	stacji ładowania EV w Polsce}
\author{}
\date{}

\begin{document}
	\maketitle
	\vspace{-1cm}
	\noindent\hrulefill
	\vspace{0.8cm}
	
	Ten dokument tłumaczy w prostym języku, skąd bierze się każda liczba i
	każdy kolor widoczny na interaktywnej mapie (\texttt{mapa\_popytu\_ev.html}).
	Nie zakłada żadnej wcześniejszej wiedzy o projekcie --- jeśli szukasz
	pełnej, technicznej dokumentacji ze wszystkimi wzorami i źródłami danych,
	zajrzyj do \texttt{dokumentacja.pdf}.
	
	\section{Co w ogóle pokazuje ta mapa?}
	
	Mapa odpowiada na jedno pytanie: \textbf{gdzie w Polsce najbardziej
		opłacałoby się postawić nową stację ładowania samochodów elektrycznych?}
	
	Każde miejsce na mapie (stacja paliw, węzeł drogowy, miejsce obsługi
	podróżnych) dostało jedną liczbę --- \textbf{szacowaną roczną energię}
	(w kWh), jaką mogłaby tam zużyć nowa stacja ładowania. Im wyższa liczba,
	tym bardziej obiecująca lokalizacja. Kolor i wielkość znaczników na
	mapie odzwierciedlają właśnie tę liczbę.
	
	\textbf{Ważne od razu na wstępie:} to nie jest pomiar, tylko
	\emph{szacunek} --- oparty na publicznie dostępnych danych o ruchu
	drogowym, liczbie samochodów elektrycznych i istniejącej konkurencji.
	Żadna z tych stacji faktycznie jeszcze nie powstała, więc nikt nie
	zmierzył tam prawdziwego zapotrzebowania.
	
	\section{Skąd wzięły się dane wejściowe?}
	
	Zanim cokolwiek policzyliśmy, zebraliśmy informacje z kilku publicznie
	dostępnych źródeł:
	
	\begin{itemize}[leftmargin=1.2em]
		\item \textbf{Ile aut przejeżdża daną drogą} --- z oficjalnych pomiarów
		ruchu drogowego (Generalny Pomiar Ruchu, GDDKiA)
		\item \textbf{Ile jest już samochodów elektrycznych i gdzie} --- z
		rejestru pojazdów (CEPiK), w podziale na powiaty
		\item \textbf{Gdzie już są stacje ładowania} --- z oficjalnego rejestru
		infrastruktury (EIPA/UDT) oraz map internetowych (OpenStreetMap)
		\item \textbf{Jaki typ zabudowy ma dany obszar} (domy jednorodzinne czy
		bloki) --- z ostatniego Spisu Powszechnego (GUS)
		\item \textbf{Ile osób faktycznie mieszka w promieniu kilometra od
			danego miejsca} --- z siatki gęstości zaludnienia GUS (kwadraty
		1$\times$1 km), żeby nie zakładać, że mieszkańcy powiatu są rozłożeni
		równomiernie
	\end{itemize}
	
	Wszystkie te informacje połączyliśmy w jedną, dużą tabelę --- 16\,232
	potencjalnych miejsc na stację, każde opisane dziesiątkami cech.
	
	\section{Jak z tych danych powstaje jedna liczba?}
	
	Tu jest sedno sprawy. Wyobraź sobie to jak przepis kulinarny: \textbf{te
		same zasady (,,przepis'') stosujemy do każdego z 16\,232 miejsc
		(,,składników'') z osobna}, żeby dostać 16\,232 osobnych wyników.
	
	Zasady różnią się w zależności od typu miejsca --- rozróżniamy dwa
	rodzaje lokalizacji:
	
	\subsection{Miejsca przy trasach (,,segment korytarzowy'')}
	
	To MOP-y, węzły drogowe i stacje paliw leżące blisko głównych szlaków
	komunikacyjnych. Tu liczy się przede wszystkim \textbf{ruch
		przejeżdżających aut}:
	
	\begin{infobox}
		Im więcej samochodów (w tym elektrycznych) przejeżdża obok, tym większy
		potencjał --- ale tylko część z nich faktycznie się zatrzyma i naładuje,
		i tylko część z tych, którzy się zatrzymują, wybierze akurat to miejsce,
		a nie konkurencję w pobliżu. Dodatkowo miejsca przy trasach z
		ponadprzeciętnym ruchem samochodów dostawczych (kurierzy, flota firmowa)
		dostają niewielką premię --- to osobna grupa kierowców, o innym wzorcu
		ładowania niż zwykli podróżni.
	\end{infobox}
	
	\subsection{Miejsca w miastach (,,segment docelowy'')}
	
	To stacje paliw w miastach i na przedmieściach. Tu liczy się przede
	wszystkim \textbf{ile elektryków mieszka w okolicy}:
	
	\begin{infobox}
		Im więcej samochodów elektrycznych jest zarejestrowanych w danym
		powiecie, tym większy lokalny potencjał --- ale ta pula nie jest już
		dzielona po równo między wszystkie miejsca w powiecie. Miejsce leżące w
		gęściej zaludnionej okolicy dostaje większy udział niż identyczne
		miejsce na rzadko zaludnionych przedmieściach, bo sprawdzamy prawdziwą
		gęstość zaludnienia w promieniu kilometra. Znaczenie ma też, czy
		mieszkańcy mają własne miejsce do ładowania w domu (właściciele domów
		jednorodzinnych rzadziej korzystają z publicznych stacji niż mieszkańcy
		bloków, którzy najczęściej nie mają własnego przyłącza).
	\end{infobox}
	
	\subsection{Co obniża lub podnosi wynik}
	
	Niezależnie od typu miejsca, na wynik wpływa też:
	
	\begin{itemize}[leftmargin=1.2em]
		\item \textbf{Konkurencja w pobliżu} --- jeśli w promieniu 2 km jest już
		dużo stacji ładowania (zwłaszcza mocnych), potencjał nowej stacji spada
		--- choć nie liniowo: dodatkowa, dziesiąta konkurencyjna stacja obok
		obniża wynik dużo mniej, niż zrobiłaby to druga czy trzecia, bo dobre
		lokalizacje z natury przyciągają wielu operatorów naraz, więc sama
		gęsta konkurencja nie powinna dyskwalifikować miejsca, które w innych
		aspektach jest bardzo dobre
		\item \textbf{Wiarygodność danych} --- jeśli dane o ruchu drogowym w
		danym miejscu są mniej pewne (bo pomiar był daleko), wynik jest lekko
		obniżany, żeby nie sugerować fałszywej precyzji
		\item \textbf{Wymogi unijne} --- lokalizacje wypełniające lukę w
		obowiązkowym rozstawie dużych stacji na głównych trasach (wymóg
		unijnego rozporządzenia AFIR) dostają dodatkową premię, tym większą, im
		ważniejsza (,,bazowa'') jest dana trasa w unijnej sieci TEN-T
	\end{itemize}
	
	\section{Co oznaczają poszczególne warstwy na mapie?}
	
	\begin{description}[leftmargin=0pt, style=nextline]
		\item[Mapa cieplna] Rozmyta, kolorowa ,,chmura'' pokazująca z lotu ptaka,
		gdzie w Polsce w ogóle warto szukać dobrych lokalizacji. Kolor pokazuje
		\emph{względną} pozycję danego miejsca (czy jest w lepszej, czy gorszej
		połowie wszystkich kandydatów) --- nie dokładną liczbę kWh. Granatowy i
		niebieski oznaczają niższy potencjał, żółty i pomarańczowy --- średni,
		czerwony --- najwyższy.
		
		\item[Zielone kółka --- najlepsze miejsca przy trasach] Wielkość i
		intensywność koloru kółka odzwierciedlają \emph{rzeczywisty} wynik ---
		większe i ciemniejsze kółko oznacza wyższy szacunek.
		
		\item[Różowo-fioletowe kółka --- najlepsze miejsca w miastach] To samo
		co powyżej, dla segmentu docelowego. W obu grupach najlepsze 200
		lokalizacji ma dodatkowo widoczny numer pozycji w rankingu (biała
		cyferka na ciemnym tle).
		
		\item[Szare kropki --- infrastruktura, która już istnieje] To miejsca,
		gdzie stacja ładowania już stoi. Nie są to rekomendacje budowy nowej
		stacji --- pokazują tylko, gdzie już jest konkurencja, i (pomocniczo) jak
		model ocenia miejsca, które realni operatorzy już wybrali.
		
		\item[Tło całych powiatów] Pokazuje, które regiony Polski mają w ogóle
		duży, niewykorzystany potencjał rynkowy --- zanim jeszcze zdecydujesz,
		które konkretne miejsce w tym regionie wybrać.
	\end{description}
	
	\subsection{Jak dokładnie liczona jest wartość w każdej warstwie?}
	
	To ważne pytanie, bo różne warstwy \textbf{nie} liczą swojej wartości w
	ten sam sposób --- część z nich mnoży kilka liczb przez siebie, jedna
	dodaje wartości wielu miejsc razem, a jeszcze inna w ogóle nie liczy
	nowej liczby, tylko porządkuje istniejące wyniki w kolejkę.
	
	\begin{description}[leftmargin=0pt, style=nextline]
		\item[Kółka (zielone i różowo-fioletowe) --- wartość to mnożenie]
		Każde kółko ma wartość policzoną jako iloczyn kilku liczb przez siebie:
		
		\begin{center}
			\small
			$\text{wynik} = \text{popyt bazowy} \times \text{pewność danych} \times \text{premia AFIR} \times \text{mnożnik ruchu} \times \text{korekta konkurencji}$
			
			(dla miejsc przy trasach dodatkowo $\times\ \text{mnożnik ruchu dostawczego}$)
		\end{center}
		
		Każdy z tych mnożników to liczba bliska 1,0 (np. 0,85--1,15) --- więc
		żaden pojedynczy czynnik nie potrafi sam z siebie drastycznie zmienić
		wyniku, ale wszystkie razem, pomnożone przez siebie, mogą znacząco
		przesunąć końcową wartość w górę lub w dół.
		
		\item[Mapa cieplna --- wartość to pozycja w kolejce, nie nowa liczba]
		Zamiast liczyć cokolwiek od nowa, mapa cieplna bierze gotowe wartości z
		kółek, \textbf{sortuje} wszystkich 6\,560 kandydatów od najsłabszego do
		najlepszego i sprawdza, na którym miejscu w tej kolejce znajduje się
		dana lokalizacja. Kolor pokazuje więc ,,czy to miejsce jest w lepszej
		czy gorszej połowie stawki'', a nie samą wartość w kWh.
		
		\item[Tło powiatów --- wartość to suma (dodawanie)] Tu dzieje się coś
		innego: bierzemy wartości wszystkich kandydatów docelowych leżących w
		danym powiecie i dodajemy je do siebie:
		
		\begin{center}
			\small
			$\text{wartość powiatu} = \text{kandydat}_1 + \text{kandydat}_2 + \text{kandydat}_3 + \ldots$
		\end{center}
		
		Duży powiat z wieloma przeciętnymi kandydatami może więc wypaść wysoko,
		nawet jeśli żaden pojedynczy kandydat w nim nie jest wybitny --- liczy
		się suma potencjału całego regionu, nie pojedyncze miejsce. Dopiero ta
		suma jest później sortowana (tak jak w mapie cieplnej), żeby dostać
		kolor.
		
		\item[Szare kropki --- ta sama wartość co kółka] Istniejące stacje mają
		wartość policzoną dokładnie tym samym mnożeniem co kandydaci pod nową
		inwestycję --- to pozwala porównać, jak model oceniłby miejsce, które
		ktoś już wcześniej wybrał.
	\end{description}
	
	\section{Ograniczenia mapy}
	
	Żeby uczciwie korzystać z tej mapy, warto pamiętać o kilku
	ograniczeniach:
	
	\begin{infobox}
		\textbf{To ranking, nie prognoza finansowa.} Liczby w kWh/rok są dobre
		do \emph{porównywania} miejsc między sobą (które jest lepsze od
		którego), ale nie należy traktować ich jako precyzyjnej, gwarantowanej
		prognozy przychodów.
		
		\textbf{Model bazuje wyłącznie na danych publicznych.} Nie uwzględnia
		np. dostępności odpowiedniego przyłącza elektrycznego, ceny gruntu,
		czy planów innych inwestorów --- tego typu informacje nie są publicznie
		dostępne.
		
		\textbf{Niektóre dane źródłowe są niepełne.} Np. flota EV w części
		powiatów (te wokół dużych miast) jest w oficjalnych rejestrach lekko
		zaniżona --- wiemy o tym i oznaczamy takie miejsca na mapie, ale nie
		naprawiamy tego w pełni, bo nie mamy dokładniejszych danych.
		
		\textbf{45 dużych miast liczone jest łącznie z okolicznym powiatem.} Dla
		obszarów jak ,,Kraków + krakowski'' dane o flocie EV nie dawały się
		rozdzielić na miasto i okoliczny powiat osobno, więc oba są traktowane
		jako jeden, połączony rynek --- co jest rozsądnym przybliżeniem, ale
		oznacza nieco mniejszą precyzję w największych aglomeracjach.
	\end{infobox}

	
\end{document}